% Source: https://artofproblemsolving.com/community/c3h2582378_discriminant_for_quadratic
% Stable identity: aops-topic-2582378-q01
For a quadratic equation $x^2+px+q=0$ with real coefficients, the sign of the discriminant $p^2-4q$ tells whether the roots are real. This problem shows that the test does not work when the coefficients are not real.
\begin{enumerate}
\item[(a)] Show that $x^2+ix+2i=0$ has no real solution.
\item[(b)] Show that $x^2+2ix-2=0$ has a positive discriminant but no real solution.
\item[(c)] Show that $x^2-(1+i)x+i=0$ has a real solution although its discriminant is not even a real number.
\end{enumerate}
